\title{Group Sequence Policy Optimization}

\begin{abstract}
We present Group Sequence Policy Optimization (GSPO), a reinforcement learning framework specifically designed for efficient, robust, and high-performing training of large language models. Distinct from earlier approaches which use token-level importance ratios, GSPO formulates importance ratios over entire sequences and executes sequence-level clipping, reward allocation, and optimization procedures. Our empirical results show that GSPO outperforms the GRPO algorithm both in training efficiency and final model performance, effectively enhances the stability of Mixture-of-Experts (MoE) RL training, and could streamline RL infrastructure design. GSPO's advantages have directly contributed to significant advancements observed in the most recent Qwen3 models.
\end{abstract}

\section{Introduction}
\label{sec:intro}

Reinforcement learning (RL) has become a central framework for advancing the capabilities of language models \citep{o1,dpsk-r1,qwq32b,qwen3}. When applied on a large scale, RL enables these models to address complex challenges, including high-level mathematical and programming tasks, by engaging in more intricate and extended reasoning.

Achieving effective scaling of RL through increased computational resources fundamentally depends on preserving consistent and resilient training behavior. Yet, modern leading RL approaches such as GRPO \citep{grpo} encounter pronounced instability during the training of very large language models, with frequent occurrences of catastrophic and unrecoverable model failures \citep{qwen3, minimax-m1}. Such instability presents a significant barrier to further advancing the performance of language models via extended RL-based training.

In this study, we demonstrate that GRPO's instability arises primarily from a core flaw: the incorrect use and breakdown of importance sampling weights within its algorithm. This flaw leads to significant high-variance noise during training, which intensifies as response sequences grow longer. The issue is exacerbated by the clipping mechanism, resulting in eventual collapse of the model.

In response to these fundamental challenges, we introduce \textbf{Group Sequence Policy Optimization (GSPO)}, an RL method designed for the training of large language models. GSPO's principal contribution is its rigorous, theory-based formulation of the importance ratio using sequence likelihood \citep{click}, which is consistent with the foundational concept of importance sampling. Furthermore, GSPO evaluates normalized rewards as the advantages among several responses to the same query, thereby strengthening the consistency between rewards assigned at the sequence level and the overall optimization objective.

Through our empirical studies, we establish that GSPO notably surpasses GRPO in terms of training stability, efficiency, and overall performance. Importantly, GSPO inherently addresses the instability commonly encountered in reinforcement learning for large Mixture-of-Experts (MoE) architectures, thereby obviating the requirement for sophisticated stabilization techniques and pointing towards a simplification of the RL workflow. These advantages have directly led to remarkable gains in the latest Qwen3 models. We anticipate that GSPO will serve as a reliable and scalable algorithmic framework, facilitating further progress in large-scale reinforcement learning with language models.

\section{Preliminaries}

\paragraph{Notation}
Throughout this paper, we represent an autoregressive language model with parameters $\theta$ as a policy $\pi_\theta$. The variable $x$ signifies a query, while $\mathcal{D}$ stands for the set of all queries considered. For a given response $y$ corresponding to the query $x$, its probability under the policy $\pi_\theta$ is represented as $\pi_\theta (y | x)=\prod_{t=1}^{|y|} \pi_\theta (y_t | x, y_{<t} )$, with $|y|$ indicating the token count in $y$. Each query-response pair $(x, y)$ receives a reward from a verifier $r$, where $r(x, y)$ maps to the interval $[ 0, 1 ]$.

\paragraph{Proximal Policy Optimization (PPO)} Leveraging trajectories collected under the previous policy $\pi_{\theta_\text{old}}$, PPO \citep{ppo} restricts the change to the policy within a local neighborhood of the prior policy by means of a clipping strategy. In particular, PPO utilizes the following policy optimization objective (here, we omit the KL regularization term for simplicity, since it is not central to our discussion):

\begin{align}
\mathcal{J}_\text{PPO}(\theta) = \mathbb{E}_{ x \sim \mathcal{D},\, y \sim \pi_{\theta_\text{old}}( \cdot | x) }
\left[ \frac{1}{|y|} \sum_{t=1}^{|y|} 
\min \left( w_{t}(\theta) \widehat{A}_{t}, \; \mathrm{clip} \left( w_{t}(\theta), 1 - {\varepsilon}, 1 + {\varepsilon}\right) \widehat{A}_{t} \right)
\right],
\end{align}

Here, the importance ratio for the token $y_t$ is given by $ w_{t}(\theta) = \frac{ \pi_{\theta} (y_{t} | x, y_{<t}) }{ \pi_{\theta_\text{old}} (y_{t} | x,y_{<t})} $, where the advantage $\widehat{A}_{t}$ of $y_t$ is derived using a separate value estimator, and $\varepsilon$ denotes the clipping threshold applied to the importance ratios.

A central practical limitation of PPO arises from its significant dependence on the value model. In most implementations, the value model’s scale closely matches that of the policy model, which significantly increases both memory consumption and computational load. Additionally, the success of the algorithm is contingent on the accuracy of its value predictions. Although obtaining a dependable value model is intrinsically difficult, the task becomes even more demanding when the aim is to support longer sequences or more sophisticated tasks at scale.

\paragraph{Group Relative Policy Optimization (GRPO)} GRPO \citep{grpo} eliminates reliance on a value function by estimating the relative advantage among a set of responses to an identical query. More precisely, GRPO aims to maximize the following objective:

\begin{align}
\label{equ:grpo}
\mathcal{J}_\text{GRPO}(\theta) = \mathbb{E}_{ x \sim \mathcal{D},\, \{y_i\}_{i=1}^G \sim \pi_{\theta_\text{old}}( \cdot | x) }
\left[ \frac{1}{G} \sum_{i=1}^{G} \frac{1}{|y_i|} \sum_{t=1}^{|y_i|} 
\min \bigg( w_{i,t}(\theta) \widehat{A}_{i,t},\, \mathrm{clip} \big( w_{i,t}(\theta), 1 - {\varepsilon}, 1 + {\varepsilon}\big) \widehat{A}_{i,t} \bigg)
\right],
\end{align}

where $G$ is the number of generated responses to each query $x$ (i.e., the group size), and the importance ratio $w_{i,t}(\theta)$ and advantage $\widehat{A}_{i,t}$ of token $y_{i,t}$ are:

\begin{align}
    w_{i,t}(\theta)=\frac{ \pi_{\theta} (y_{i,t} | x, y_{i,<t}) }{ \pi_{\theta_\text{old}} (y_{i,t} | x,y_{i,<t})},\quad\ 
    \widehat{A}_{i,t} = \widehat{A}_{i} = \frac{r(x, y_i) - \mathrm{mean} \left( \{ r(x, y_i) \}_{i=1}^G \right) }{ \mathrm{std} \left( \{ r(x, y_i) \}_{i=1}^G \right) },
\end{align}

respectively, where all the tokens in $y_i$ share the same advantage as $\widehat{A}_{i}$.

\section{Motivation}
\label{sec:motivation}

As models become larger, more sparse (as seen in Mixture-of-Experts architectures), and generate longer responses, deploying a substantial rollout batch size becomes essential for optimal hardware performance in RL settings. To enhance sample efficiency, it is commonly adopted to subdivide this sizeable batch of rollouts into several smaller mini-batches when computing gradient updates. This, however, inherently leads to an off-policy learning scenario: here, responses $y$ are drawn from an older version of the policy, $\pi_{\theta_\text{old}}$, rather than from the currently optimized policy $\pi_\theta$. Such a context clarifies the rationale behind employing the clipping techniques in PPO and GRPO, which serve to exclude samples that deviate too far from the present policy, thereby ensuring reliable gradient estimates.

Although approaches such as clipping attempt to address the off-policy discrepancy, our analysis reveals a deeper problem within GRPO: \textit{the objective itself is not well-defined}. This issue is especially pronounced when employing large-scale models for tasks involving extended responses, frequently resulting in severe model degradation. The core of this ill-posedness originates from an incorrect use of importance sampling weights. Importance sampling is fundamentally intended to approximate the expected value of a function $f$ with respect to a target distribution $\pi_\text{tar}$, by appropriately re-weighting data sampled from a behavioral distribution $\pi_\text{beh}$:

\begin{align}
\mathbb{E}_{z \sim \pi_\text{tar}} \left[ f(z) \right] = \mathbb{E}_{z \sim \pi_\text{beh}} \left[ \frac{ \pi_\text{tar} (z) }{ \pi_\text{beh} (z)} \, f(z) \right].
\end{align}

Importantly, this approach depends on taking an average across numerous samples ($N \gg 1$) drawn from the behavior distribution $\pi_\text{beh}$, so that the importance weight $\frac{ \pi_\text{tar} (z) }{ \pi_\text{beh} (z)}$ can appropriately adjust for discrepancies between distributions.

On the other hand, GRPO incorporates the importance weighting factor $\frac{ \pi_{\theta} (y_{i,t} | x, y_{i,<t}) }{ \pi_{\theta_\text{old}} (y_{i,t} | x,y_{i,<t})}$ at every token step $t$. This factor relies on a single sampled token $y_{i,t}$ from the subsequent-token distribution $\pi_{\theta_\text{old}}( \cdot | x, y_{i, <t})$ and, therefore, is inadequate in fulfilling the correction objective between distributions. Rather than correcting, it injects significant high-variance noise into the computed training gradients. Such variance tends to build up through lengthy sequences, with the effect magnified further by the clipping method employed. Our empirical studies reveal that this process can trigger a form of model collapse that proves difficult or impossible to recover from. After collapse, attempts to resume training—including restoring earlier model checkpoints, exhaustive hyperparameter adjustments (such as those for clipping bounds), lengthening generated sequences, or altering RL queries—have not succeeded in reversing the degeneration.

The preceding analysis highlights a central flaw in the architecture of GRPO. The ineffectiveness of token-level importance weighting underscores an essential principle: \textbf{the optimization objective's granularity must align with that of the reward}. Because rewards are assigned to complete sequences, attempting to implement off-policy correction for individual tokens proves to be unsuitable. Consequently, we are led to abandon token-level objectives in favor of employing importance weights and conducting optimization at the \textit{sequence level} instead.

\section{Algorithm}

\subsection{GSPO: Group Sequence Policy Optimization}

Although the token-level importance weight $\frac{ \pi_{\theta} (y_{i,t} | x, y_{i,<t}) }{ \pi_{\theta_\text{old}} (y_{i,t} | x,y_{i,<t})}$ poses issues within GRPO, we note that the \textit{sequence-level} importance weight $\frac{ \pi_{\theta} (y | x) }{ \pi_{\theta_\text{old}} (y | x)}$ holds clear theoretical significance in language generation tasks. Specifically, it quantifies the degree of divergence between the response $y$ generated under $\pi_{\theta_\text{old}} (\cdot | x)$ and that under $\pi_{\theta} (\cdot | x)$. This aligns inherently with sequence-level rewards and offers a principled basis for evaluating or guiding the clipping strategy.

Based on this straightforward observation, we propose the \textbf{Group Sequence Policy Optimization (GSPO)} algorithm. GSPO employs the following sequence-level optimization objective:

\begin{align}
\mathcal{J}_\text{GSPO} (\theta) =
\mathbb{E}_{ x \sim \mathcal{D},\, \{y_i\}_{i=1}^G \sim \pi_{\theta_\text{old}}( \cdot | x) }
\left[ 
\frac{1}{G} \sum_{i=1}^{G}
\min \bigg( s_{i}(\theta)  \widehat{A}_{i},\,
\mathrm{clip} \big( s_{i}(\theta),\, 1 - {\varepsilon},\, 1 + {\varepsilon} \big) \widehat{A}_{i} \bigg)
\right],
\label{equ:gspo}
\end{align}

where we adopt the group-based advantage estimation:

\begin{align}
\widehat{A}_{i} = \frac{r(x, y_i) - \mathrm{mean} \left( \{ r(x, y_i) \}_{i=1}^G \right) }{ \mathrm{std} \left( \{ r(x, y_i) \}_{i=1}^G \right) },
\end{align}

and define the importance ratio $s_{i}(\theta)$ based on sequence likelihood \citep{click}:

\begin{align}
s_{i}(\theta) = \left( \frac{ \pi_{\theta} (y_i | x) }{ \pi_{\theta_\text{old}} (y_i | x)} \right)^{\frac{1}{|y_i|}}
=
\exp \left( \frac{1}{|y_i|} \sum_{t=1}^{|y_i|} \log \frac{ \pi_{\theta} (y_{i,t} | x, y_{i,<t}) }{ \pi_{\theta_\text{old}} (y_{i,t} | x,y_{i,<t})} \right).
\end{align}

Consequently, GSPO implements clipping at the level of whole responses rather than at the token level, thereby filtering out excessively ``off-policy'' samples from the gradient computation and ensuring consistency with sequence-level reward assignment and optimization objectives. Additionally, we employ length normalization in $s_{i}(\theta)$ to mitigate variance and maintain $s_{i}(\theta)$ within a standardized numerical interval. Without such normalization, minor variations in the likelihoods of certain tokens could cause significant swings in the sequence-level importance ratio, and this would necessitate different clipping bounds depending on response length. Furthermore, it should be emphasized that the clipping intervals used by GSPO generally differ by orders of magnitude from those of prior methods (e.g., GRPO), primarily because importance ratios are defined differently across these algorithms.

\subsection{Gradient Analysis}
\label{subsec:gradient}

We can derive the gradient of the GSPO objective as follows (clipping is omitted for brevity):

\begin{align}
\nabla_{\theta} \mathcal{J}_\text{GSPO} (\theta)
=&\ 
\nabla_{\theta} \mathbb{E}_{ x \sim \mathcal{D},\, \{y_i\}_{i=1}^G \sim \pi_{\theta_\text{old}}( \cdot | x) }
\left[ 
\frac{1}{G} \sum_{i=1}^{G} s_{i}(\theta) \widehat{A}_{i}
\right] \\
= &\ 
\mathbb{E}_{ x \sim \mathcal{D},\, \{y_i\}_{i=1}^G \sim \pi_{\theta_\text{old}}( \cdot | x) }
\left[ 
\frac{1}{G} \sum_{i=1}^{G}
s_{i}(\theta) \widehat{A}_{i}
\, \nabla_{\theta} \log s_{i}(\theta)
\right] \\
=&\ 
\mathbb{E}_{ x \sim \mathcal{D},\, \{y_i\}_{i=1}^G \sim \pi_{\theta_\text{old}}( \cdot | x) }
\left[ 
\frac{1}{G} \sum_{i=1}^{G} 
\left( \frac{ \pi_{\theta} (y_i | x) }{ \pi_{\theta_\text{old}} (y_i | x)} \right)^{\frac{1}{|y_i|}} \widehat{A}_{i} 
\, \frac{1}{|y_i|} \sum_{t=1}^{|y_i|} \nabla_{\theta} \log \pi_{\theta} (y_{i,t} | x, y_{i,<t}) 
\right].
\label{equ:gspo_grad}
\end{align}

For comparison, the gradient of the GRPO objective is as follows (note that $\widehat{A}_{i,t} = \widehat{A}_{i}$):

\begin{align}
\nabla_{\theta} \mathcal{J}_\text{GRPO} (\theta) 
=&\ 
\nabla_{\theta} \mathbb{E}_{ x \sim \mathcal{D},\, \{y_i\}_{i=1}^G \sim \pi_{\theta_\text{old}}( \cdot | x) }
\left[ \frac{1}{G} \sum_{i=1}^{G} \frac{1}{|y_i|} \sum_{t=1}^{|y_i|} 
w_{i,t}(\theta) \widehat{A}_{i,t}
\right] \\
=&\
\mathbb{E}_{ x \sim \mathcal{D},\, \{y_i\}_{i=1}^G \sim \pi_{\theta_\text{old}}( \cdot | x) }
\left[ \frac{1}{G} \sum_{i=1}^{G} \widehat{A}_{i} 
\times \left(\frac{1}{|y_i|} \sum_{t=1}^{|y_i|} 
\frac{ \pi_{\theta} (y_{i,t} | x, y_{i,<t}) }{ \pi_{\theta_\text{old}} (y_{i,t} | x,y_{i,<t})} 
\nabla_{\theta} \log \pi_{\theta} (y_{i,t} | x, y_{i,<t}) \right)
\right].
\end{align}

The primary difference between GSPO and GRPO centers on \textit{the manner in which each method assigns weights to the gradients of token log likelihoods}. GRPO applies a token-level ``importance weight'' given by $\frac{ \pi_{\theta} (y_{i,t} | x, y_{i,<t}) }{ \pi_{\theta_\text{old}} (y_{i,t} | x,y_{i,<t})}$ to each token. These resulting weights are inherently uneven and can range within $(0, 1+\varepsilon]$ when $\widehat{A}_{i} > 0$, or $[1-\varepsilon, +\infty)$ if $\widehat{A}_{i} < 0$, and their nontrivial magnitude can aggregate during training, causing uncertain or unstable learning behavior. Conversely, GSPO adopts uniform weighting for every token within a response, thereby circumventing the instability introduced by GRPO's varied token-level weighting.

\subsection{GSPO-token: A Token-level Objective Variant}

In contexts such as multi-turn RL, there may be a need for more granular advantage modification beyond the sequence level. To address this, we present a token-level modification of GSPO, referred to as \textbf{GSPO-token}, which enables the assignment of advantages on a per-token basis:

\begin{align}
\mathcal{J}_\text{GSPO-token}(\theta)
=
\mathbb{E}_{ x \sim \mathcal{D},\, \{y_i\}_{i=1}^G \sim \pi_{\theta_\text{old}}( \cdot | x) }
\left[
\frac{1}{G} \sum_{i=1}^G \frac{1}{|y_i|} \sum_{t=1}^{|y_i|}
\min\Bigg(
s_{i,t}(\theta) \widehat{A}_{i,t},\,
\mathrm{clip}\big(s_{i,t}(\theta), 1 - {\varepsilon}, 1 + {\varepsilon}\big) \widehat{A}_{i,t}
\Bigg)
\right],
\label{equ:gspo-token}
\end{align}

where

\begin{align}
s_{i,t}(\theta) 
= \mathrm{sg} \left[ s_{i}(\theta) \right]  \cdot \frac{ \pi_{\theta} (y_{i,t} | x, y_{i,<t}) }{ \mathrm{sg} \left[ \pi_{\theta} (y_{i,t} | x, y_{i,<t}) \right] },
\end{align}

and $\mathrm{sg}[\cdot]$ denotes only taking the numerical value but stopping the gradient, corresponding to the \texttt{detach} operation in PyTorch. The gradient of GSPO-token can be derived as:

\begin{align}
\nabla_{\theta} \mathcal{J}_\text{GSPO-token}(\theta)
=&\ 
\nabla_{\theta} \mathbb{E}_{ x \sim \mathcal{D},\, \{y_i\}_{i=1}^G \sim \pi_{\theta_\text{old}}( \cdot | x) }
\left[ 
\frac{1}{G} \sum_{i=1}^{G}
\frac{1}{|y_i|} \sum_{t=1}^{|y_i|} 
s_{i,t}(\theta) \widehat{A}_{i,t}
\right] \\
=&\ 
\mathbb{E}_{ x \sim \mathcal{D},\, \{y_i\}_{i=1}^G \sim \pi_{\theta_\text{old}}( \cdot | x) }
\left[
\frac{1}{G} \sum_{i=1}^{G} s_i (\theta) \cdot \frac{1}{|y_i|} \sum_{t=1}^{|y_i|} 
\widehat{A}_{i,t} \frac{ \nabla_{\theta} \pi_{\theta} (y_{i,t} | x, y_{i,<t}) }{ \pi_{\theta} (y_{i,t} | x, y_{i,<t}) }
\right] \\
=&\ 
\mathbb{E}_{ x \sim \mathcal{D},\, \{y_i\}_{i=1}^G \sim \pi_{\theta_\text{old}}( \cdot | x) }
\left[ 
\frac{1}{G} \sum_{i=1}^{G}  \left( \frac{ \pi_{\theta} (y_i | x) }{ \pi_{\theta_\text{old}} (y_i | x)} \right)^{\frac{1}{|y_i|}}
\cdot \frac{1}{|y_i|} \sum_{t=1}^{|y_i|}
\widehat{A}_{i,t} \nabla_{\theta} \log \pi_{\theta} (y_{i,t} | x, y_{i,<t})  
\right].
\label{equ:gspo-token_grad}
\end{align}

It should be observed that the expression $\frac{ \pi_{\theta} (y_{i,t} | x, y_{i,<t}) }{ \mathrm{sg} \left[ \pi_{\theta} (y_{i,t} | x, y_{i,<t}) \right] }$ evaluates numerically to 1. Consequently, $s_{i,t}(\theta)$ yields the same numerical result as $s_{i}(\theta)$. Upon examination of Equation~\eqref{equ:gspo} alongside~\eqref{equ:gspo-token}, as well as Equation~\eqref{equ:gspo_grad} with~\eqref{equ:gspo-token_grad}, one finds that GSPO-token and GSPO are numerically equivalent in terms of their optimization objectives, clipping criteria, and theoretical gradients, provided the advantages assigned to each token within the response $y_i$ are held constant (that is, $\widehat{A}_{i,t} = \widehat{A}_{i}$ for all $t$). Nevertheless, GSPO-token additionally allows for the advantage value to be tailored individually for each token, offering increased flexibility.

\section{Experiments and Discussion}

\subsection{Empirical Results}

We utilize a cold-start model that has been fine-tuned from Qwen3-30B-A3B-Base, and present both the training reward progression and the model’s performance on the AIME'24 benchmark (calculating average Pass@1 from 32 samples), LiveCodeBench (202410-202502, computing average Pass@1 over 8 samples), and CodeForces (measured using Elo Rating). For reinforcement learning training, each rollout batch is subdivided into four mini-batches to perform gradient updates. In implementing GSPO, the clipping intervals for Equation~\eqref{equ:gspo} are fixed at 3e-4 for the lower bound and 4e-4 for the upper bound. For comparison, we use GRPO as a baseline approach and assign clipping bounds in Equation~\eqref{equ:grpo} to 0.2 and 0.27, respectively, after thorough parameter tuning to achieve comparability. It is important to highlight that, for effective convergence in MoE RL using GRPO, the Routing Replay training method is required—a topic elaborated in \S~\ref{subsec:routing_replay}. By contrast, \textit{GSPO eliminates the necessity for this method}.

As illustrated in Figure~\ref{fig:results}, GSPO enables stable progress during training. Notably, \textbf{GSPO achieves ongoing improvements in performance by scaling up training compute, periodically refreshing the query set, and lengthening generation output}. In addition, GSPO outperforms GRPO in terms of training efficiency, reaching higher training accuracy and superior benchmark results given equivalent computational resources and query budgets. Finally, GSPO has been effectively utilized for RL training in the newest Qwen3 models, providing strong evidence for its capability to harness the benefits of RL scaling in large language models.

\subsection{Curious Observation on Clipping Fractions}

An important difference between GSPO and GRPO lies in their approach to clipping: GSPO clips responses at the sequence level, whereas GRPO applies clipping at the token level. As depicted in Figure~\ref{fig:clipping}, this distinction results in a difference of nearly two orders of magnitude in the proportion of clipped tokens across the two methods, and this large gap persists even when the clipping intervals are varied. Remarkably, GSPO achieves better training efficiency than GRPO, even though it clips a considerably higher number of tokens, leading to fewer tokens being available for training or for gradient calculation. This seemingly paradoxical observation—that training is more efficient when a much greater share of tokens are excluded—suggests that token-wise gradient estimates in GRPO are fundamentally noisy and not well-suited for effective sample exploitation. Conversely, GSPO’s strategy of generating gradients at the response level delivers a clearer and more efficient training signal.

\subsection{Benefit of GSPO for MoE Training}
\label{subsec:routing_replay}

\paragraph{Background}
While Reinforcement Learning (RL) training of dense neural architectures is relatively stable, the sparse activation property inherent to MoE models presents distinct stability issues. Notably, when utilizing the GRPO algorithm, the MoE models display pronounced \textit{expert-activation volatility}, which hampers the convergence of RL optimization. Specifically, even a small number of gradient updates can lead to substantial shifts in the set of experts selected for producing identical responses. For instance, experiments with the 48-layer Qwen3-30B-A3B-Base model reveal that following each RL gradient update, approximately 10\% of the experts activated by the new policy $\pi_{\theta}$ for a fixed rollout diverge from those chosen by the preceding policy $\pi_{\theta_\text{old}}$. This effect becomes especially acute in deeper MoE architectures, leading to severe oscillations in the token-level importance weights $w_{i,t}(\theta) = \frac{ \pi_{\theta} (y_{i,t} | x, y_{i,<t}) }{ \pi_{\theta_\text{old}} (y_{i,t} | x,y_{i,<t})}$, as further elaborated in \S~\ref{sec:motivation} and~\ref{subsec:gradient}. As a result, the instability in these weights impedes proper convergence during RL training.

\paragraph{Our Previous Approach} In our earlier work, we addressed this issue by introducing the \textbf{Routing Replay} training method. This approach involves recording the experts selected by $\pi_{\theta_\text{old}}$, then applying these identical routing patterns in $\pi_{\theta}$ during the calculation of importance ratios $w_{i,t}(\theta) = \frac{ \pi_{\theta} (y_{i,t} | x, y_{i,<t}) }{ \pi_{\theta_\text{old}} (y_{i,t} | x,y_{i,<t})}$. By doing so, the two probability models, $\pi_{\theta} (y_{i,t} | x, y_{i,<t})$ and $\pi_{\theta_\text{old}} (y_{i,t} | x,y_{i,<t})$, employ the same set of network experts for each token $y_{i,t}$. This alignment helps to maintain consistent importance ratios at the token level and supports stable optimization on the activated subnetwork across successive gradient steps. As depicted in Figure~\ref{fig:routing_replay}, Routing Replay proves to be a crucial mechanism to guarantee stable convergence of the GRPO algorithm for Mixture-of-Experts (MoE) models.

\paragraph{Benefit of GSPO} While Routing Replay enables MoE models to achieve convergence under GRPO training, its reliance on preserving and reusing routing configurations introduces extra communication and memory costs and can inhibit the effective utilization of the model's overall capacity. GSPO, by contrast, as illustrated in Figure~\ref{fig:results}, completely removes the requirement for Routing Replay. It straightforwardly computes the importance ratios $s_i(\theta)$ in the standard manner, ensuring stable optimization and normal convergence. The main observation is that GSPO is concerned exclusively with the likelihood over full sequences (i.e., $\pi_{\theta} (y_i | x)$) rather than with the likelihood at the token level (i.e., $\pi_{\theta} (y_{i,t} | x, y_{i,<t})$). Given that MoE models consistently retain their ability for language modeling, the likelihood of entire sequences does not experience significant variability. In essence, GSPO addresses the instability in expert activation typical of MoE models at its source, removing the necessity for elaborate techniques such as Routing Replay. As a result, the approach not only streamlines and fortifies the training dynamics but also permits the model to exploit its total capacity without imposed restrictions.

\subsection{Benefit of GSPO for RL Infrastructure}

Due to differences in numerical precision between training engines (such as Megatron) and inference engines (such as SGLang and vLLM), it is standard practice to recalculate the likelihoods of generated responses under the previous policy $\pi_{\theta_\text{old}}$ using the training engine. Notably, GSPO relies on sequence-level likelihoods for its optimization process, rather than focusing on token-level measures; as a result, sequence-level approaches inherently exhibit greater robustness to such precision mismatches. This feature of GSPO allows for the optimization process to utilize likelihood estimates obtained directly from the inference engine, eliminating the need for additional recomputation via the training engine. Such capability is particularly advantageous in contexts like partial rollout, multi-turn RL, and within training-inference disaggregated frameworks.

\section{Conclusion}

We introduce Group Sequence Policy Optimization (GSPO), a novel reinforcement learning framework tailored for optimizing large language models. Anchored in the central concepts of importance sampling, GSPO utilizes sequence likelihoods to define importance ratios, and implements clipping, reward assignment, and updates at the sequence level. When benchmarked against GRPO, GSPO achieves markedly better stability during training, heightened efficiency, and enhanced performance. Its effectiveness is especially pronounced in large-scale reinforcement learning scenarios with MoE models, underpinning the significant advancements observed in the recent Qwen3 models. With GSPO serving as a robust and scalable foundation, we are committed to expanding the limits of reinforcement learning, anticipating consequential progress in artificial intelligence.

\end{document}